% 习惯.tex
% Auto-generated from 感悟论语.odt by tools/odt_to_tex.py

\setpartplant{artwork/peach.png}
  {桃夭 · 《诗·周南》}
  {性相近也，习相远也。}
  {《论语·阳货》}
  {习惯}
\part{习惯}

\chapter{行为决定习惯}

甘地说过，“信念决定思想，思想决定语言，语言决定行为，行为决定习惯，习惯决定性格，性格决定命运”。整个链条的表述连贯流畅，似乎只需有一个“信念”作为起点，随后便可顺其自然地养成“习惯”、塑造“性格”、收获成功的或至少是有意义的“命运”。但事实上，“信念”的确立与“思想”的形成往往需要机缘，并非通过艰辛的苦求而一定会有所得。相比之下，“语言”容易被组织成文，以形成对自我有益的心理暗示，但其功效也仅限于此。接下来的“行为”则是最为困难的一环。因为决定一个人最终能够到达何种层次与境界的关键不在于他的想法有多么理性、系统、深刻，而取决于他在每一天精神状态最佳、大约8小时的工作时间里究竟能以什么样的品质做着什么样的事，这件事能否如实而完美地体现他最初的信念，以及他能够将这样的工作状态日复一日地延续多久。以这样的笃行作为根基，之后的“习惯”与“性格”便会水到渠成。而最终的“命运”则非人的意志与努力所能决定，能够“知天命、顺天命”便好。

基于这样的想法，我们也就不难理解为何儒家重视行动与实践且将其作为好学的标准了。例如，

\begin{classic}
《论语·学而篇第一 1.6》 子曰：「弟子入则孝，出则弟，谨而信，泛爱众，而亲仁。行有余力，则以学文。」

《论语·公冶长篇第五 5.14》 子路有闻，未之能行，唯恐有闻。

《论语·雍也篇第六 6.3》 哀公问：「弟子孰为好学？」孔子对曰：「有颜回者好学，不迁怒，不贰过。不幸短命死矣，今也则亡，未闻好学者也。」

《论语·子路篇第十三 13.1》 子路问「政」。子曰：「先之，劳之。」请益。曰：「无倦。」
\end{classic}

\entrydate{2022.03.14}

\chapter{好习惯的养成}

早起可以说是人人皆知的好习惯。试想在清晨时分，第一缕阳光透过树丛映照于窗棂，轻快的鸟鸣引领着世界从静谧中缓缓走出，在这样清静、安详的气氛中静心冥想，或是在灵感的激励下奋笔疾书，真的是修炼与提升自我的大好时机。可虽然益处显而易见，但人人却又知实践早起之难。由此可见，想要有个如此简单的好习惯还真不是一件容易的事。

如果想要养成好习惯，去除已有的坏习惯，就得让自己从那种舒服的环境中跳出来。只有觉得不舒服、不习惯、紧张而有压力，才会让自己变得越来越出色和游刃有余。好习惯的养成，一定要用清醒的头脑去审视自己每一天的计划与实际行动。即便是百忙于琐事之中，也要善于抓住蕴含于其中的诸多时间片断，以便时时调整自己的心态、精神状态和实时目标。这时不妨审视一下自己的心态是否宁静？精神状态是否积极？目标是否足够明确而不至无所适从或者因陷于细节的泥淖而因小失大？如此日积月累，才会逐渐形成习惯乃至产生人生境界的质变。

对于成人而言，其养成新的好习惯的过程与孩童不同。小孩子有大人为他们提出的明确要求和时时监督，且没有生活中种种琐事的牵纠而使其分心。因此，他们无需太多的自控力，只需照章做事，即可逐渐将习惯养成。而成人则要面临或抵抗多方面的干扰与影响，难得有大片的闲暇时间可以专注于自己的内心从而展开深入的反省与思考，亦不再有来自“权威”的监督与警告，多年的积习也让人更容易囿于舒适圈内而不思改变与进取。因此，成人只要稍不留意或懈怠，便会放松对自己的管控，不形成新的坏习惯就已经不错了，好习惯的养成更是无从谈起。

由此看来，只有永不倦怠，对思想和行为时刻保持警醒地审视，才能从混乱中跳出来，逐渐养成更多的好习惯。这正如曾子所言，“战战兢兢，如临深渊，如履薄冰”。希望在众多好习惯的推动下，我们的人生变得更加丰富，具有积极的意义！

\begin{classic}
《论语·泰伯篇第八 8.3》 曾子有疾，召门弟子曰：「启予足！启予手！诗云：『战战兢兢，如临深渊，如履薄冰。』而今而后，吾知免夫！小子！」
\end{classic}

\entrydate{2017.04.11}

\chapter{循序渐进养成好习惯}

好的生活和学习习惯的养成，处事原则的固化于心，以及恰当的待人接物方式的练就，都不是一蹴而就的，也绝对不是将所有的经验、规律、原则都理性地想清楚并条理清晰地整理出来形成标准和规范后就能够立刻体现于行动的。它们需要伴随着个人意志力与执行力缓慢、稳步的提升而逐渐增加砝码。如同学习不宜贪多，若将期望的好习惯一下子实践得过多，当因个人暂且较差的执行力和较低的完成率而不能看到与之对应的立竿见影的效果时，反而会打击信心，容易因灰心丧气而一曝十寒甚至半途而废。这也就是《论语·公冶长篇第五 5.14》所说的：

\begin{classic}
子路有闻，未之能行，唯恐有闻。
\end{classic}

所以正确的方法是从点滴小事做起，从简单的任务做起，而且不能贪多。当看到自己能够很好地完成既定目标时，情绪得到鼓舞，精神得到激励，然后再逐渐地增加每日需要完成的习惯性任务的种类、数量和强度，由此便形成了良性循环，不断提升自我，臻于完美。

\begin{classic}
《论语·公冶长篇第五 5.14》 子路有闻，未之能行，唯恐有闻。
\end{classic}

\entrydate{2018.08.31}

\chapter{养成坚持写日记的习惯}

自2018年1月22日起，我与部门的两位同事立下了“好习惯坚持三个月”的承诺。于是，我的日记写得比去年更加频繁了，基本上每天都能够码出一些文字来。可以说到目前为止，写日记成为我每天早起洗漱后的第一件事，业已固化为一个生活习惯了——这确是令人欣慰的。

虽然我既不是什么名人，没有什么惊天动地的事情值得去写；亦与作家相去甚远，锤炼不出什么优美隽永的文字，但即便如此，写日记的习惯所能带来的好处仍是不言而喻的。比如，日记是对生活中主要事件、活动、心理、思想的忠实记录。在将来回首往事的时候，它能够栩栩如生地展现和还原出曾经的心路历程。若以此作为材料和根据，就有利于明辨事理和考量思想、态度、行为上的得失，从而有益于“如切如磋、如琢如磨”般地自我提升和改进。这也正是曾子所说的“吾日三省吾身——为人谋而不忠乎？与朋友交而不信乎？传不习乎？” 。再者，写日记其实也是一个熟能生巧的过程。通过这样的自我训练，理性思维与结构化提炼、组织素材的能力会得到不断巩固和强化，对文字的驾驭也会逐渐得心应手。这些能力的塑造不仅仅有助于个人的有效学习，更是将输入型学习升华为输出型学习从而对他人产生积极影响的有力武器。

日记除了能够对自己的人生产生助推作用外，若将其整理出来与大家分享，也说不定会给他人带来意料之外的启发和帮助。即便这些文字本身并不务实地传播知识和技能，也不急功近利地贩卖管理学、成功学之类的心灵鸡汤，可能就是那么几行自然而朴实的讲述和描绘，就能够令读者觉得“原来也有相似的人在这里呀”。这就像是我曾经在豆瓣上读到“跑跑小羊”写的那篇《你不曾真的离去，你永远在我心里……》中的一句话：

\begin{classic}
不经意的遇见，触动了心灵，产生出久远的共鸣和激励，自己在纯粹理性世界中的求索也不再觉得孤单。
\end{classic}

再退一步讲，就算日记中的文字和写日记这个活动本身没有实质性的意义，但通过每天的笔耕不辍，仍可以于无形之中养成重视承诺并努力践行的习惯，从而有效地提高了意志力与执行力。以如此“更好的自我”再去训练和养成个人生活中其他方面的好习惯，不也是一件颇有收获和成就感的事吗？

\begin{classic}
《论语·学而篇第一 1.15》 子贡曰：「贫而无谄，富而无骄，何如？」子曰：「可也，未若贫而乐，富而好礼者也。」子贡曰：「诗云：『如切如磋，如琢如磨』，其斯之谓与？」子曰：「赐也，始可与言诗已矣，告诸往而知来者。」

《论语·学而篇第一 1.4》 曾子曰：「吾日三省吾身，为人谋而不忠乎？与朋友交而不信乎？传不习乎？」
\end{classic}

\entrydate{2018.06.27}

\chapter{五维度反省}

今日清晨读到傅佩荣微信公众号的文章《每天睡前十分钟培养一种习惯》。其中提到：

\begin{classic}
每天10分钟做什么？自我反省。今天发生了什么事？从这些事情，我得到的教训是什么？每天10分钟，养成这种习惯。一个月之后，你看自己就已经开始透明化了。
\end{classic}

以我自己的实践相对照，目前在手写笔记本上能够大致每天用钢笔写的“回顾与小结”便属于这样的反省。其内容于“家庭”、“做事”、“学习”、“开发”、“习惯”共五个维度展开。其中，“家庭”关乎孝悌与关爱\footnote{《论语·学而篇第一 1.2》 有子曰：「其为人也孝弟，而好犯上者，鲜矣；不好犯上，而好作乱者，未之有也。君子务本，本立而道生。孝弟也者，其为仁之本与！」}，是对内的。“做事”是对外的，要反省为他人做事时是否坚守原则、问心无愧、尽心尽责\footnote{《论语·学而篇第一 1.4》 曾子曰：「吾日三省吾身，为人谋而不忠乎？与朋友交而不信乎？传不习乎？」}。“学习”则是为了自我修炼，在追求君子的永无止境的道路上不断提升\footnote{《论语·子张篇第十九 19.6》 子夏曰：「博学而笃志，切问而近思，仁在其中矣。」}。“开发”即为我的数值算法软件开发工作，这是践行“格物致知”、“博学近思”的具体方式，并在“君子上达”精神的引领下争取做出一番功业\footnote{《论语·宪问篇第十四 14.23》 子曰：「君子上达，小人下达。」}。“习惯”既帮助我在最大程度上维持稳定、健康、长久的工作状态，也在日积月累中逐渐培养出个人品质与风格。

在纸面上完成上述五个维度的反省工作，往往10分钟还不够。相信将这个方法一直保持下去，学问可以逐渐变得深厚、笃实，研究风格与工作态度变得愈加纯粹，在算法开发上兼顾整体架构与代码细节，形成与理论框架的完美呼应。

\entrydate{2022.07.16}

\chapter{写纸质日记的益处}

用纸笔写日记的方式为我们提供了一个没有空间与时间限制、可以自由抒写心思感悟的天地。日记既可以工整美观、精致考究，亦可以字迹潦草、恣意洒脱。有时，文字像意识流般地娓娓道来；有时，页面里充满的则是零落而无意义的涂鸦。

身处充斥着碎片化信息的世界，面对人事纷扰令心灵难有片刻安宁的情境，我们之所以仍要坚持写日记，就是因为我们总是不可避免地在人生的某个阶段会对今后的走向感到些许迷茫，会对当下所投入的事业产生一丝怀疑，对于能否收获期待的果实与成功亦没有十足的信心。有所思、有所感，便会有倾吐心思的欲望。

然而，指望通过与好友的交流来获得某种慰藉与启示并非总能成功。毕竟，人与人之间心灵的契合需要某种难以捉摸的机缘，朋友间保持发人深省的沟通也需要彼此共同有志于学和不断的自我提升。对于这一点，孔夫子早就说过：“可与共学，未可与适道；可与适道，未可与立；可与立，未可与权。”所以，此时若将所有的迷茫、怀疑、气馁诉诸纸笔，则会有助于思维的明澈、精神的提振、注意力的汇聚与信心的恢复。通过在文字中的自我反省，最终实现精神与心理上的提升。而这种提升往往能够有效地隔离来自外界的干扰，帮助我们形成独立的思想与意志。

此外，《乾卦·文言》中提到“修辞立其诚”。从中便可看出，修饰言辞是促进我们锤炼出纯正与真诚品格的重要手段。而写日记正好提供了充分的训练机会，让我们能够用精确、严谨、优雅、美好的语言阐明我们至善的心意、坚定的决心与远大的志向。由此，我们的行动也才会变得愈加笃实。而这一点恰如圣雄甘地所说：

\begin{classic}
Always aim at complete harmony of thought and word and deed. Always aim at purifying your thoughts and everything will be well.
\end{classic}

\entrydate{2020.10.05}

\chapter{背单词一年总结}

我从本科毕业后，由于既没有了英语考级的压力，也不得不将大部分时间用于科研，所以直到博士毕业的这段时间内，就没有专门在英语学习上投入精力。同时也一直觉得，自己整天都在读英文原版的数学、物理、计算机书籍，怎么着应该算是捎带着学了外语吧！之后从学校毕业又走上工作岗位，我也仍旧觉得作为一名普通的科技工作者，只要准确掌握了完整的英语语法并具有足够的专业词汇量，应付学习和工作绝对是绰绰有余了。这样一来，十四年的时光里，我的英语水平除了熟练度有所改善、词汇量略有零星的增加外，便一直没有实质性的进步。这正可谓是“画地自限”了。

本以为自己就这样不思进取、“自暴自弃”下去，幸好去年跟爱人一起系统地复习了英语语法，便借此机缘重新激活了坚持学习英语的热情与毅力（虽有佳肴 ，弗食 ，不知其旨也 ；虽有至道 ，弗学 ，不知其善也 。是故学然后知不足 ，教然后知困 。知不足 ，然后能自反也 ；知困 ，然后能自强也 。故曰 ：教学相长也 。——《礼记·学记第十八》）。自2017年11月13日起，我开始在欧路词典中背单词，到今天正好整整一年。一开始只是背学习和研究中查过的生词。后来，自2018年1月29日始，进一步学习Merriam Webster应用中的“每日一词”。这里出现的单词往往非常生僻，还会将词源（拉丁语、法语等）和演变的来龙去脉讲得一清二楚。所以，初学的时候极不适应。从2018年5月13日始，摘录词典中的例句和书中精辟、优美的句子，目的是为生词的学习提供一个语境，既便于记忆，也能学到地道、专业、生动、恰如其分的表达方法。

由此收获的改进与体会总结如下：

\begin{itemize}
\item 熟能生巧：通过每天的练习，英语单词发音的连贯性和准确度得到了提高。
\item 迎难而上：对于阅读或学习过程中一下子出现大量生词的情况不再抵触和恐惧，能够不厌其烦地逐个查清楚。
\item 温故知新：掌握了许多已学单词的另一重含义，印证了“温故而知新”的道理。
\item 学而时习：充分利用欧路词典提供的卡片学习和批量重学功能，定期调出已经学过的单词重新学习。只有通过不断地清零和反复学习，才能抵抗持续的遗忘。
\item 细致笃实：发现一个规律就是基本上所有单词都具有不止一种词性。例如，一个词可以同时做名词和动词用，或者名词和形容词用。由此深刻地体会到，对于我们自以为熟悉的单词所做的望文生义十有八九都是靠不住的或者是有所偏差的。要想真正掌握英语单词，就必须摒弃学生时代为了应试而采用的根据上下文“猜测”的方法——做阅读理解必用的招数——那着实是投机取巧、浮躁功利的坏毛病。真正的学习要求我们绝不能放过阅读过程中遇到的每一个生词，必须老老实实地查词典，且应像是学习数学概念一样，明确它们的词性、含义、程度、适用范围。
\item 专业准确：口语虽然听起来地道，且极易流行，但是也更容易被时间淘汰。相比之下，书面语也许是因为被一帮学究们薪火相传，所以词义更为准确、到位，留存得也更为稳定、持久。这样看来，学习一定量的专业甚至是生僻词汇是非常必要的。
\end{itemize}

\begin{classic}
《论语·雍也篇第六 6.12》 冉求曰：「非不说子之道，力不足也。」子曰：「力不足者，中道而废。今女画。」

《论语·为政篇第二 2.11》 子曰：「温故而知新，可以为师矣。」

《论语·学而篇第一 1.1》 子曰：「学而时习之，不亦说乎？有朋自远方来，不亦乐乎？人不知而不愠，不亦君子乎？」
\end{classic}

\entrydate{2018.11.12}

\chapter{远离手机、找回灵魂}

智能手机及其应用导致分心的问题想必每个人都或多或少地有所意识和感触。不过，本文在开篇便要强调的观点是：即便在空闲或是休息的时候使用手机，同样会影响到正常的学习与工作。

这是因为当我们面对手机时，大脑始终被无关紧要的碎片化信息干扰，从而无法得到片刻的休息与安宁，仿佛有一群挥之不去的蚊蝇萦绕在身边，令人心烦意乱。在这样的环境与心境下，又怎么可能进行严谨、缜密、深入的理性思考？更何况，对于一般人而言——其实也不只是一般人了，许多大师级的人物也如此——像是体育竞技赛前运动员的热身，进入理性思考的状态往往需要一个没有外界干扰、安静舒缓、自由闲适的情思过程，比如，午后的散步、清晨的冥想。在经历这个过程之后，才有可能形成思想的火花，并酝酿成科学与艺术的灵感。《论语·八佾3.23》有云：

\begin{classic}
子语鲁大师乐，曰：“乐其可知也：始作，翕如也；从之，纯如也，皦如也，绎如也，以成。”
\end{classic}

由此可以看出，理性思考如同一阙乐章绝不只是最终得到的那个结论，而是一个完整的过程，甚至可能还相当的漫长。试想，若在我们思考的乐章中浑入来自各种媒体的噪音，哪怕极为短暂，乐则不能成，思亦无所得。所以可以准确地说，多种社交、音频与视频媒体所破坏的并不是思考本身，而是形成思考的过程。想来这也就是对于当下许多人来说，慢阅读、深刻思考、理性分析变得尤为困难的原因之一吧！在还没有思考之前，他们就已经被干扰了，使得真正的思考根本无法开始。对此不可不慎之又慎！

以上所说的是各种媒体对于思考过程的影响。那么对于媒体本身所传达的信息，我们是否又深入地理解与消化了呢？也未必。因为大多数时候，我们对于那些被转发多次的好文——无论是时间管理、自我提升，还是文艺鉴赏、科技知识——也并不是全神贯注阅读的，更谈不上做笔记、写心得了。而至于根据自身情况将其付诸实践，持久坚持以至于行成良好习惯，则是鲜有所闻。相比之下，回想我们在手机和社交应用出现之前是如何面对媒体和信息的：

仅就我自己来说，十几前年我听广播的时候，基本上不做什么其他无关的事情，就是专心致志地聆听。所以直至今日，仍旧清晰记得从电影《简爱》录音剪辑中所听出的人物的音容笑貌，对白中透出的淡淡忧伤与浓浓爱意；心中也还依然存留着一汪澄澈而《宁静的湖泊》——而那是电台播放并解说的英国作曲家柯兹的作品。许久以前的读书，也一定是正襟危坐，带着一种神圣感在世界名著中徜徉，其间不会被外界干扰或打断。而那个时候聆听音乐，除了必备的音乐知识尚不足备外，真可以算得上是“欣赏”了。往往是一盘磁带翻来复去地听了不知多少遍，分分秒秒都是全身心地投入。正因为如此，古典音乐中经典、华丽、优雅的音符能够多年以来一直驻留在自己的心里——这就像是电影《肖申克的救赎》里的主人公安迪，在他的脑海与心中一直吟唱的那段咏叹调（《费加罗的婚礼》苏珊娜与伯爵夫人的合唱）。

鉴于过往，想必我们可以有所省悟和收获了。至少在力所能及的情况下，尽量远离手机，以避免无论做什么事情还未进入全神贯注的状态时，便被手机夺走了灵魂。如果这种情形得不到改观，一切的行为不过是在走过场而不走心，最终难有所得。

\entrydate{2017.02.13}
